\section*{Limitations}
\phantomsection\label{sec:limitations}

The measures employed in this study are deliberately shallow. Orthographic
similarity compares trigram inventories and is therefore sensitive to the
transcription conventions of a particular translation and not solely to the
language that translation transcribes. Phonetic similarity depends upon Double
Metaphone, an encoder developed for English, whose treatment of Philippine
phonology is approximate, uniform across all sixteen languages, and liable to
compress genuine contrasts.

Both measures are computed over Biblical translation, which constitutes a single
register produced by translators working to their own conventions, and eight of
the sixteen corpora cover only the New Testament (\cref{tab:seeds}). Register
and translator effects are consequently confounded with language effects
throughout, and no comparison reported here is able to separate them.

The analysis presented in \cref{tab:metric_comparison} addresses the most
tractable of these concerns, namely corpus size, and establishes that the
disparity in scale between the two measures is not attributable to it. It does
not address the remainder. Those residual concerns bear upon the magnitudes
reported in \cref{sec:synthesis} rather than upon the ordering of language
relationships, and it is upon the ordering that the convergent findings of that
section principally rest.
